\scanpage{73}
\margin{Si $\dim A=0$ ou $1$, il n'y a \unsure{rien à prouver}, \unsure{on peut supposer|supposons} $\dim A\geq2$}
Soit ${\unsure{\mathscr H|\mathcal H}}^1$ le \deleted{\unsure{gr.}} \unsure{ens.} des classes de fibrés principaux homogènes (\unsure{triv} par extension fid. plat quasi compacte de la base). On aura $H^1(X-a,G)\subset\mathscr H^1(X-a,G)$, qui lui est \unsure{partie} d'un foncteur coh. en $G$.

On a si \deleted{$\unsure{\dim A=0}$} $G$ est \deleted{\unsure{une variété abélienne}} \inserted{\uline{propre} sur $k$}
\[
H^1(X-a,G)\deleted{=\unsure{H^1(X,G)}}=0\qquad(\text{loc trivial}=\text{trivial}),
\]
\unsure{donc on prouve|donc on prouvera} $\unsure{\forall|\text{tt}}$ élément \unsure{de|d'un} $H^1(X-a,G)$ \unsure{provient} de $\mathscr H^1(X-a,L)$, $L$ le ss-gpe linéaire \inserted{\uline{connexe}}.

Or pour un gpe linéaire \inserted{connexe}, on sait que \deleted{$H^1=\mathscr H^1$} $H^1(G)=\mathscr H^1(G)$, donc on \deleted{\unsure{est}} aura
\[
H^1(X-a,G)=\deleted{H^1(X-a,L)}\operatorname{Im}\bigl(H^1(X-a,L)\longrightarrow H^1(X-a,G)\bigr)
\]
\[
\operatorname{Im}\bigl(\mathscr H^1(X-a,L)\longrightarrow\mathscr H^1(X-a,G)\bigr)
=H^1(X-a,L)/\operatorname{Im}H^0(X-a,G/L)
\]
Or $H^0(X-a,G/L)=H^0(X,G/L)$ qui \unsure{s'envoie} dans $H^1(X,L)=0$, donc $\uline{\ensuremath{H^1(X-a,G)}}=\uline{\ensuremath{H^1(X-a,L)}}$.

D'ailleurs, comme $X$ est factoriel, $H^1(X-a,G_m)=0$ d'où $H^1(X-a,L)=H^1(X-a,U)$, $U$ le \unsure{gpe} unipotent. Cette dernière se \unsure{dévisse} par les $H^1(X-a,G_a)$, qui sont nuls si $\dim A>2$, mais \unsure{non} \unsure{pour} \deleted{\unsure{si}} $\dim A=2$.

\[
\text{d)}\quad H_A^{\unsure{3}}(G)=H^2(X-a,G)
\]
De façon générale, on a \deleted{$H^i(X-a,G)=$} une suite exacte de faisceaux
\[
0\longrightarrow\underline{\unsure{\mathscr O|O}}_X(G)\unsure{\mathbin{\text{---}}|\longrightarrow}\underline{\unsure{\mathscr R|R}}_X(G)\mathbin{\text{---}}\underline{\unsure{\mathscr D|D}}_X(G)\mathbin{\text{---}}0
\]
d'où, puisque $\underline{\mathscr R}_X(G)$ est flasque
\[
H^i(\underline{\mathscr O}_X(G))=H^{i-1}(\underline{\mathscr D}_X(G))\quad\text{si }i\geq2
\]
% archive sheet number: 72
